\section{Closed-form parametric constants and the assumption-surface floor}
\label{app:closedforms}

\paragraph{The partial-adjustment law.} The simplest transition in the class is
the partial-adjustment law \eqref{eq:affine}.
Its state and sanction slopes equal $1-\delta$ and $\delta$ at both ends of their
bounds, and every curvature modulus is zero. Where the lower and upper bounds on
either slope coincide, the transition is of the form \eqref{eq:affine} on the
square (\Cref{prop:bracket}). Under a fixed sanction $s$, the unique rest point
of every transition in the class is $x=s$, because \textup{(F1)--(F5)} exclude a
range of inaction.

\paragraph{Status of the constants in derivative-free form.} The derivative-free
parametric-argmax results of \citet{KlatteKummer2002} and the
set-valued analysis of \citet{AubinFrankowska1990} provide the
appropriate route from \emph{complete-objective} strong concavity and
objective-difference bounds to Lipschitz selectors. Stage curvature by
itself is not such a route. The complete target-objective modulus
$c_T$, where invoked, is the scalar form of the diagonal strict
concavity condition of \citet{Rosen1965}.

This appendix lists the closed-form constants entering
\Cref{ass:sender,ass:target,ass:concavity} and the Schur condition
\eqref{eq:Schur}. The paper-facing primitives are the discount and
transition objects $\beta_T,\beta_S\in[0,1)$ and
$\F\colon[0,1]^2\to[0,1]$ satisfying \textup{(F1)--(F5)};
the stage payoffs $\Ut,\Us\colon[0,1]^3\to\R$ (continuous and twice
continuously differentiable on $(0,1)^3$); the basin pivots
$\bar x_T,\bar x_S\in[0,1]$, derived from $\Ut,\Us$ via the
threshold characterization in \Cref{sec:mechanism:basin}; and the
first-derivative bounds $L_{U_T},L_{U_S},L_{U_a}$ used in the
sanction- and action-sensitivity estimates, defined as suprema of
$|\partial\Ut/\partial s|$, $|\partial\Us/\partial s|$, and
$|\partial U/\partial a|$ on $[0,1]^3$.

The value-function state bound has its own numerator. Let
$L_{S,x}^{\alpha}$ bound the inherited-state variation of
$\Us(x,\alpha(x),s)$, and let $\gamma_F$ be the fixed-sanction transition
modulus in \eqref{eq:F-state-lip}. In primitive applications both are
discharged from derivative and transition bounds; neither is a sanction
slope, and under the affine law $\gamma_F=1-\delta$. The target-side
analogue (a closed-loop state modulus $\gamma_T$ and the Lipschitz
constant $K_T=L_{T,x}^{\mathrm{cl}}/(1-\beta_T\gamma_T)$ of the target's invariant set in the optional
target-Lipschitz route) does not enter the existence construction for the class under
the lifting condition, whose invariant set for the target is a norm ball with a lower-slope floor and
no Lipschitz clause. It is instantiated on that
class after the fact: with $\gamma_T$ the sharp loop constant
$g_{\mathrm{hi}}+d_{\mathrm{hi}}L_{\mathrm{resp}}/\mu$ and
$L_U=\max\{v_{\mathrm{hi}},v_{\mathrm{hi}}+n_{\mathrm{hi}}\}+cL_{\mathrm{resp}}/\mu$
($v+n+cL_{\mathrm{resp}}/\mu$ on the quadratic--linear family), the equilibrium target value is
$K_T=L_U/(1-\beta_T\gamma_T)$-Lipschitz whenever $\beta_T\gamma_T<1$
(\path{VTstar_lipschitz}, \path{Examples/PlacableSenderTargetLipschitz.lean}).

\subsection{Derived and separately verified constants}

The following list distinguishes algebraically derived bounds from
positive curvature or selector moduli that must be verified where
used:
\begin{itemize}[leftmargin=2em]
\item $K_S=L_{S,x}^{\alpha}/(1-\beta_S\gamma_F)$, closed-form Lipschitz
constant of the sender's invariant set (derived from $\Vs\in\mathrm{Lip}_{K_S}$
by Banach iteration under $\beta_S\gamma_F<1$; the recursion is \citet[Thm.~4.2(c)]{Hinderer2005} and its inheritance by the invariant set is Corollary~1 to Theorem~3.2 of \citet{StokeyLucas1989}).
\item $c_\sigma=\inf_{[0,1]^3}\bigl(-\partial^2\Us/\partial s^2\bigr)>0$,
the stage strong-concavity modulus imposed in
\Cref{ass:concavity}. It controls the static sender problem only.
\item No general dynamic-selector modulus $L_\sigma$ is derived from
$c_\sigma$. Whenever an existence theorem uses a Lipschitz dynamic
sender selector, its named analytic package must instead verify strong
concavity of the complete objective $J_S$ and the relevant
parametric objective-difference bound
(\citealp{KlatteKummer2002,DoraszelskiPakes2007}).
\item $K_T^{\mathrm{IFT}},K_S^{\mathrm{IFT}}\ge0$ objective-difference
moduli of the joint target objective in the value-pair coordinates: a
sup-norm perturbation of size $\varepsilon$ in $\Vt$ (respectively
$\Vs$) changes the objective's variation between two actions by at most
$K^{\mathrm{IFT}}\varepsilon$ times the action distance. This is the
objective-difference input of
\path{argmax_lipschitz_of_strongConcave_lipschitzDiff}, verified by
each analytic package that uses it. It is not a finite-dimensional
Hessian bound, and smoothness of the stage payoff on a neighborhood of
the cube does not supply it: a uniformly strongly concave objective can
have a maximizer that is not Lipschitz in sup-norm perturbations of the
objective.
\item $c_T>0$ joint-target complete-objective strong-concavity modulus
in $a$, verified as part of the leader-regularity package wherever it
is used; continuity of second derivatives alone does not imply its
positivity.
\item $L_{\alpha,\Vt}=K_T^{\mathrm{IFT}}/c_T$,
$L_{\alpha,\Vs}=K_S^{\mathrm{IFT}}/c_T$, Implicit-Function Lipschitz
constants of the target equilibrium policy in the value-pair
coordinates. The canonical derivation from C${}^2$ target stage
payoffs plus an Arrow-Hurwicz-Uzawa constraint qualification (without
an interiority hypothesis on the optimal compliance choice) is
\citet{RinconZapateroSantos2009}, which proves continuous
differentiability of the value function and uniqueness of the KKT
multipliers; the IFT gradient strong-monotonicity in derivative-free form taken
as a primitive here is
the gradient-monotonicity content of this uniqueness statement. The
encyclopedic reference for the parametric-argmax Lipschitz stability
results from which this derivative-free form descends is
\citet[Ch.~2 + Ch.~7]{KlatteKummer2002}.
\end{itemize}

Every parametric constant above is either (a) algebraically derived
from verified moduli via a closed-form expression
(e.g.\ $K_S$, $L_{\alpha,V_*}$), or (b) a bound verified from
the complete objective by the analytic package that uses it (the
curvature moduli $c_\sigma$, $c_T$ and the objective-difference moduli
$K_T^{\mathrm{IFT}}$, $K_S^{\mathrm{IFT}}$); where such a bound is read
off a smooth stage payoff, the derivative bound must be verified on the
closed cube, since differentiability on the open cube does not bound
curvature at the boundary.
The positive curvature moduli $c_\sigma$ and $c_T$ must be verified
from the corresponding stage or complete objective; compactness only
guarantees that a continuous curvature bound attains its infimum, not
that the infimum is positive. The optional moduli
$L_{S,x}^{\alpha},\gamma_F$ must
themselves be discharged from the specified payoff, transition, and
selector bounds; they are not silently identified with sanction
derivatives. The structural
assumptions \Cref{ass:sender,ass:target,ass:concavity} are
qualitative; they enter the paper-facing surface as conditions on
the payoff functions and the basin, not as numerical primitives.

\begin{center}
\small
\textbf{Table 2.} Derived and separately verified intermediates.\\[2pt]
\begin{tabular}{@{}p{0.20\linewidth}p{0.73\linewidth}@{}}
\toprule
\textbf{Symbol} & \textbf{Closed form / interpretation}\\
\midrule
$K_S$ & $L_{S,x}^{\alpha}/(1-\beta_S\gamma_F)$, Lipschitz constant of the Sender's
invariant set under \textup{(F6-Q)} and $\beta_S\gamma_F<1$; under the
affine law, $\gamma_F=1-\delta$.\\
$L_\sigma$ & Optional dynamic Sender-selector modulus; available only
from complete-objective strong concavity plus a parametric
objective-difference bound, not from A3 alone.\\
$K_T^{\mathrm{IFT}},K_S^{\mathrm{IFT}}$ & Cross-Hessian bounds on
the joint Target integrand.\\
$L_{\alpha,\Vt},L_{\alpha,\Vs}$ & $K_T^{\mathrm{IFT}}/c_T$,
$K_S^{\mathrm{IFT}}/c_T$, IFT Lipschitz of $\alpha$ in
$(\Vt,\Vs)$.\\
$L_{TS},L_{ST}$ & Row-Lipschitz off-diagonals of the joint Bellman
operator (\Cref{sec:existence}).\\
$R_T,R_S$ & Sup-norm bounds on the balls of the invariant sets; the
smallest invariant radii are $\|\Ut\|_\infty/(1-\beta_T)$ and
$\|\Us\|_\infty/(1-\beta_S)$ (\Cref{app:jointiteration}).\\
$\bar x_T,\bar x_S$ & High-basin tipping thresholds: infima of the
flip sets \eqref{eq:FT-def}--\eqref{eq:FS-def}
(\Cref{prop:tipping}); $\bar x=\max(\bar x_T,\bar x_S)$.\\
$x_S^0$ & Static Sender flip ($m=0$ level); contested band
$[x_S^0,\bar x_S)$ (\S\ref{sec:flip}).\\
$\bar x_T^*,\bar x_S^*$ & Low-basin ceilings: suprema of the dual
flip sets (\S\ref{sec:mechanism:bistability});
$\bar x_*=\min(\bar x_T^*,\bar x_S^*)$.\\
$d_{\F}(x),\eta_{\F}(\bar x)$ & Full-sanction drift
$d_{\F}(x)=\F(x,1)-x$ and its positive minimum on $[0,\bar x]$.\\
$\widehat N_{\F}(\bar x)$ & Explicit general-transition bound
$\lceil\bar x/\eta_{\F}(\bar x)\rceil$ on the worst-case shock-entry
duration under \textup{(F1)--(F4)} and full-sanction ascent on $[0,\bar x]$.\\
$\eta_\delta(\bar x),\widehat N_\delta(\bar x)$ &
$\delta(1-\bar x)$ and
$\lceil\bar x/[\delta(1-\bar x)]\rceil$, the affine specialization of
the general minimum-drift bound. The sharper affine choice is
$N_*(\bar x,\delta)=\lceil\log(1-\bar x)/\log(1-\delta)\rceil$
(\Cref{thm:shock,cor:shock-transition-strength}).\\
\bottomrule
\end{tabular}
\end{center}

The notation choice $R_T, R_S$ for the sup-norm bounds of the invariant sets (in
place of the older $M_T, M_S$) avoids collision with $M=(\alpha^M,
\sigma^M)$, the myopic equilibrium
(\Cref{def:myopic}).

\subsection{Certified comparative-static windows}

The full-adjustment comparative statics at the corner economies hold on
exact parameter intervals rather than only at isolated calibration
points. \Cref{tab:compstat-windows} records the intervals and the
constraint that binds at each edge.

\begin{table}[ht]
\centering
\small
\begin{tabular}{@{}p{0.17\linewidth}p{0.14\linewidth}p{0.61\linewidth}@{}}
\toprule
\textbf{Parameter} & \textbf{Certified window} & \textbf{What binds at the endpoints}\\
\midrule
\multicolumn{3}{@{}l@{}}{\emph{Eaton--Engers economy}}\\
\addlinespace[2pt]
Sender patience $\beta_S$
& $[6/13,\ 4/7)$
& The continuation drag must cover the worst stage drop at the lower
edge and remain below the strict foresight premium at the upper edge.\\
\addlinespace[2pt]
Target patience $\beta_T$
& $[598/709,\ 1)$
& The Target entrenchment condition binds at the lower edge; the upper
edge is the discount bound.\\
\addlinespace[2pt]
Adjustment speed $\delta$
& $[3/7,\ 2/3)$
& The Sender corner pin binds below and the strict foresight premium binds above.\\
\midrule
\multicolumn{3}{@{}l@{}}{\emph{Interest-group economy}}\\
\addlinespace[2pt]
Sender patience $\beta_S$
& $[0,\ 1)$
& The whole range is certified because the defiance-corner
stage slope is already positive.\\
\addlinespace[2pt]
Target patience $\beta_T$
& $[14/17,\ 1)$
& The Target condition and the entrenchment condition bind at the lower edge.\\
\addlinespace[2pt]
Adjustment speed $\delta$
& $[7/27,\ 1)$
& The Target condition and the invariant-set membership condition bind together at
the lower edge.\\
\bottomrule
\end{tabular}
\caption{Certified parameter windows for the full-adjustment
comparative statics. One deep parameter varies in each row while the
other primitives remain at their example values. Closed endpoints
correspond to weak conditions; open endpoints to strict conditions or
the discount bound.}
\label{tab:compstat-windows}
\end{table}

At the certified equilibria the value functions also have exact state
derivatives, the shadow prices of political entrenchment. They are
$82/11$ for the Target and $16/3$ for the Sender at the
Eaton--Engers economy, and $80/11$ and $4/3$ at the interest-group
economy. Each obeys the boundary-envelope recursion
\[
\text{slope}=U_x+\beta(1-\delta)\,\text{slope}.
\]
The positive signs certify that the equilibrium incumbents value a
harder political state, while the exact magnitudes are
example-specific rather than general predictions.
